% Independently written integration of the stationary overlap calculation.
% Compatible with working_derivation.tex: amsmath,amssymb,mathrsfs,booktabs.
\providecommand{\E}{\mathbb E}
\providecommand{\dd}{\mathrm d}
\providecommand{\diag}{\operatorname{diag}}

\section{A stationary replica closure through formal order $N^{-1/2}$}
\label{sec:adopted-stationary}

The calculation in this section supplies a closed, parameter-computable
first-order approximation to the covariance energy. It uses a stationary
replica fluctuation expansion rather than a pathwise tangent expansion about
a chaotic trajectory. Its status is a formal large-$N$ closure conditional on
the assumptions stated below. It is not a uniform error theorem for networks
of size $1000$, and exact inversion of its finite matrices does not complete
the second-order theory.

The assumptions are a specified stationary mixing state, a single relevant
mean-field saddle, independent directed edge disorder in the supplied dense
ensemble, a nonsingular inter-replica fluctuation operator, and sufficient
delocalization and uniform moment bounds to justify the cavity sums and the
stationary/frequency limits. Coherent population feedback is retained before
performing the bulk expansion. Frozen networks and other attractors require
a separate state-selection rule.

\subsection{Model, normalization and independent stationary replicas}
For a cell $i$ in population $a\in\{E,I\}$,
\[
 \dot x_i=-x_i+\sum_jJ_{ij}\phi(x_j)+\sqrt N I_a+h_a^{\rm ext},
 \qquad \phi(x)=2\Phi(x),\quad r_i=\phi(x_i).
\]
Let $f_a=N_a/N$, $F=\diag(f_E,f_I)$, and write
\[
 J_{ij}=J^0_{ab}/\sqrt N+X_{ij},\qquad
 \E X_{ij}^2=G_{ab}/N,\qquad
 \E X_{ij}^3=T_{ab}/N^{3/2}.
\]
Put $W=GF$, $\Theta=TF$, and $\epsilon=N^{-1/2}$. Removing finitely many
probe cells does not change the original normalization $N$ in these formulas.
For temporal means $m_i=\langle r_i\rangle_t$ and centered rates
$\delta r_i=r_i-m_i$, define
\[
 C_{ij}=\langle\delta r_i\delta r_j\rangle_t,\quad
 \nu_N=\frac1N\operatorname{tr}C,\quad
 \mathcal E_N=\frac1N\operatorname{tr}(C^2),\quad
 d_N=\frac{\mathrm{PR}}N=\frac{\nu_N^2}{\mathcal E_N}.
\]
Take two independent draws from the same stationary measure at fixed $J$,
denoted by superscripts $1,2$. Directly expanding the square gives the exact
identity
\begin{align*}
 \mathcal E_N
 &=\frac1N\sum_{ij}
   \E_1[\delta r_i^{(1)}\delta r_j^{(1)}]
   \E_2[\delta r_i^{(2)}\delta r_j^{(2)}]\\
 &=N\E_{1,2}\left[
       \frac1N\sum_i\delta r_i^{(1)}\delta r_i^{(2)}\right]^2.
\end{align*}
Independent initial conditions may replace independent stationary draws only
if they select this same stationary measure. The two copies share the
connectivity and the static operating-point labels; their temporal
fluctuations are independent conditional on those quenched quantities.

Use $\widehat g(\omega)=\int e^{-i\omega t}g(t)\dd t$ and
$H(\omega)=(1+i\omega)^{-1}$. On a record of length $T_r$, introduce
\[
 q_a(\omega_1,\omega_2)=\frac1{N_a}\sum_{i\in a}
   \delta\widehat r_i^{(1)}(\omega_1)
   \delta\widehat r_i^{(2)}(\omega_2).
\]
The spectral energy is $N T_r^{-2}\E|f^\top q|^2$, integrated with
$\dd\omega_1\dd\omega_2/(2\pi)^2$. The calculation below suppresses
frequency arguments when they are fixed. Numerical spectra stored as masses
already include these integration weights and are summed once.

\subsection{One-site inputs and the first-order marginal correction}
The leading conditional law is
\[
 x_a(t)=\mu+u_a(t),\quad
 \mu=h_a+\sqrt{\Delta_a}z,\quad z\sim\mathcal N(0,1),\quad
 \E u_a(t+\tau)u_a(t)=K_a(\tau),\quad q_a=K_a(0).
\]
Here $q_a$ is an input variance; $q_a(\omega_1,\omega_2)$ above is a
replica overlap, distinguished by its arguments. All label averages below
refer to this leading Gaussian label measure. Put $s_a^2=1+q_a$ and
\begin{align*}
 m(\mu)&=2\Phi(\mu/s_a),\\
 \alpha(\mu)&=\partial_\mu m
 =\sqrt{\frac2{\pi s_a^2}}e^{-\mu^2/(2s_a^2)},\qquad
 \beta(\mu)=\partial_\mu\alpha=-\mu\alpha/s_a^2,\\
 c_\mu(\tau)&=\operatorname{Cov}(r(t+\tau),r(t)\mid\mu).
\end{align*}
For $k=K_a(\tau)$ the latter is the explicit Gaussian integral
\[
 c_\mu(k)=\frac2\pi\int_0^{\arcsin(k/s_a^2)}
     \exp\left[-\frac{\mu^2}{s_a^2(1+\sin\vartheta)}\right]\dd\vartheta.
\]
Its Price derivative and mixed derivative are
\[
 D_\mu(k)=\partial_kc_\mu
 =\frac{2e^{-\mu^2/(s_a^2+k)}}{\pi\sqrt{s_a^4-k^2}},\qquad
 B_{21,\mu}(k)=\tfrac12\partial_\mu D_\mu(k).
\]

Let $y=H*r=m+v$, so $v=H*\delta r$, and write
$P_\mu(\tau)=\operatorname{Cov}(v(t+\tau),v(t)\mid\mu)$. The input
third-cumulant coefficients are the following one-site functionals:
\begin{align*}
 t_{0a}&=\sum_b\Theta_{ab}\E_b[m^3],\\
 M_a(\tau)&=\sum_b\Theta_{ab}\E_b[mP_\mu(\tau)],\\
 D_a(t_1,t_2,t_3)&=\sum_b\Theta_{ab}
       \E_b[\operatorname{cum}(v(t_1),v(t_2),v(t_3)\mid\mu)].
\end{align*}
These are computable from Gaussian expectations of $\phi$ and causal leak
convolutions. For example the rate third cumulant is a trivariate Gaussian
expectation of three $2\Phi$ factors minus its disconnected terms; filtering
its three arguments gives the cumulant of $v$. It does not require an
uncomputed reservoir-response average.

Let $h_1,\Delta_1,K_1$, and $q_1=K_1(0)$ denote the marginal shifts. With
$\sigma=(\mu-h_a)/\Delta_a$, define
\[
 \delta q_\mu=q_{1a}+\sigma M_a(0),\qquad
 \delta K_\mu(\tau)=K_{1a}(\tau)+\sigma M_a(\tau),
\]
and $D_{0a}=D_a(0,0,0)$,
$D_{sa}(\tau)=[D_a(\tau,\tau,0)+D_a(-\tau,-\tau,0)]/2$.
At a fixed leading label, the conditional corrections are
\begin{align*}
 m_1&=\tfrac12\beta\,\delta q_\mu
       +\tfrac16D_{0a}\partial_\mu^3m,\\
 \alpha_1&=\tfrac12\partial_\mu^3m\,\delta q_\mu
       +\tfrac16D_{0a}\partial_\mu^4m,\\
 c_1(\mu,\tau)&=(\partial_qc_\mu)\delta q_\mu
       +D_\mu(K_a(\tau))\delta K_\mu(\tau)
       +\tfrac16D_{0a}\partial_\mu^3c_\mu\\
 &\quad+D_{0a}\alpha\beta+
      [D_{sa}(\tau)-D_{0a}]B_{21,\mu}(K_a(\tau)),\\
 \partial_qc_\mu&=\tfrac12\partial_\mu^2c_\mu+\alpha^2-D_\mu.
\end{align*}
The label measure contributes the signed first-order weight
\[
 w_a^{(1)}(z)=\frac{t_{0a}}{6\Delta_a^{3/2}}(z^3-3z)
       +\frac{h_{1a}}{\sqrt{\Delta_a}}z
       +\frac{\Delta_{1a}}{2\Delta_a}(z^2-1).
\]
It acts as a differential correction to averages, not as a positive tail
probability distribution.

The coefficients solve the linearized one-site self-consistency equations
\begin{align*}
 0&=\E_a[w_a^{(1)}m+m_1],\\
 \Delta_{1a}&=\sum_bW_{ab}\E_b[w_b^{(1)}m^2+2mm_1],\\
 K_{1a}''(\tau)&=K_{1a}(\tau)
       -\sum_bW_{ab}\E_b[w_b^{(1)}c_\mu(\tau)+c_1(\mu,\tau)],\\
 K_{1a}'(0)&=0,\qquad K_{1a}(\infty)=0,\qquad q_{1a}=K_{1a}(0).
\end{align*}
The zero first-order mean-rate shift is the asymptotic compensated-drive
constraint. Its distinction from exact finite-network balance is addressed
below. The one-site outputs needed by the overlap theory are
\begin{align*}
 A_{0a}&=\E_a\alpha^2,&
 A_{1a}&=\E_a[w_a^{(1)}\alpha^2+2\alpha\alpha_1],\\
 Q_{0a}&=\E_a[c_\mu(\omega_1)c_\mu(\omega_2)],\\
 Q_{1a}&=\E_a[w_a^{(1)}c_\mu(\omega_1)c_\mu(\omega_2)
        +c_1(\mu,\omega_1)c_\mu(\omega_2)
        +c_\mu(\omega_1)c_1(\mu,\omega_2)],\\
 \nu_0&=\sum_af_a\E_ac_\mu(0),&
 \nu_1&=\sum_af_a\E_a[w_a^{(1)}c_\mu(0)+c_1(\mu,0)].
\end{align*}
Both factors in $Q_0$ or $Q_1$ belong to the same neuron's label.

\subsection{A stationary trajectory weight and two triple overlaps}
At leading order the cross-replica covariance fluctuation of a cavity input is
$H_1H_2Wq$. A Gaussian covariance perturbation of the product of rates has
conditional response $\alpha^2$. This yields
\[
 q=\xi+\mathcal T_0q,\qquad
 \mathcal T_0=\diag(A_0)XW,\quad X=H_1H_2,
\]
where the leading covariance of the row source, with the $N/T_r^2$
normalization understood, is $\Sigma_0=\diag(Q_{0a}/f_a)$.

The third edge cumulant introduces the centered square
$s(t)=y(t)^2-\langle y^2\rangle_t$ and the population overlaps
\[
 w_a^A=\frac1{N_a}\sum_{i\in a}
        \widehat s_i^{(1)}(\omega_1)\widehat v_i^{(2)}(\omega_2),\qquad
 w_a^B=\frac1{N_a}\sum_{i\in a}
        \widehat v_i^{(1)}(\omega_1)\widehat s_i^{(2)}(\omega_2).
\]
The derivative of a squared filtered rate is
\[
 \frac{\delta y(t)^2}{\delta x(\sigma)}
       =2H(t-\sigma)y(t)\phi'(x(\sigma)).
\]
Consequently its label-conditioned stationary response is
\begin{align*}
 \Lambda_\mu(\tau)&=2\theta(\tau)e^{-\tau}
            [m\alpha+\ell_\mu(\tau)],\\
 \ell_\mu(\tau)&=\operatorname{Cov}(v(t+\tau),\phi'(x(t))\mid\mu)
                  =\tfrac12\partial_\mu(H*c_\mu)(\tau).
\end{align*}
The last identity follows by differentiating a Gaussian rate covariance with
respect to its common mean: exchange symmetry makes the two derivative terms
equal. Define
\[
 \ell_\mu^+(\omega)=\int_0^\infty
        e^{-(1+i\omega)\tau}\ell_\mu(\tau)\dd\tau,
 \qquad \widehat\Lambda_\mu(\omega)=2m\alpha H(\omega)+2\ell_\mu^+(\omega).
\]
The $\ell$ term retains a same-replica temporal correlation. Conditional
independence of the two replicas does not justify dropping it.

Write $b_a=\E_a[\alpha\beta]$ and set
\begin{align*}
 \mathcal C&=\tfrac12\diag(b)\Theta,\qquad
 \mathcal T_{\mathrm{loc}}=\diag(A_1)XW,\\
 \mathcal D_A&=\diag\{\E_a[\widehat\Lambda_\mu(\omega_1)\alpha]\}
                    H_2XW,\\
 \mathcal D_B&=\diag\{\E_a[\alpha\widehat\Lambda_\mu(\omega_2)]\}
                    H_1XW.
\end{align*}
The term $\mathcal T_{\mathrm{loc}}$ contains the one-site law correction.
The closed augmented linear system, to the required formal order, is
\begin{equation}\label{eq:adopted-augmented}
 \begin{pmatrix}
 I-\mathcal T_0-\epsilon\mathcal T_{\mathrm{loc}}&-\epsilon\mathcal C&-\epsilon\mathcal C\\
 -\mathcal D_A&I&0\\-\mathcal D_B&0&I
 \end{pmatrix}
 \begin{pmatrix}q\\w^A\\w^B\end{pmatrix}
 =\begin{pmatrix}\xi\\\zeta^A\\\zeta^B\end{pmatrix}.
\end{equation}
Each block is two-by-two in population space. The factor $1/2$ in
$\mathcal C$ is the third-cumulant Price coefficient: there are three
orderings of two inputs in one replica and one in the other, out of $3!$.

Eliminating the two triple overlaps is a Schur complement:
\begin{align*}
 q&=(I-\mathcal T_0-\epsilon\mathcal T_1)^{-1}S,\\
 \mathcal T_1&=\mathcal T_{\mathrm{loc}}+\mathcal C(\mathcal D_A+\mathcal D_B),
 \qquad S=\xi+\epsilon\mathcal C(\zeta^A+\zeta^B).
\end{align*}
For $m_{2a}=\E_a[m\alpha^2]$ and
$e_a(\omega)=\E_a[\alpha\ell_\mu^+(\omega)]$, this is explicitly
\begin{align*}
 \mathcal T_1={}&\diag(A_1)XW\\
 &+\diag(b)\Theta\{
       2\diag(m_2)X+\diag(e(\omega_1))H_2
                        +\diag(e(\omega_2))H_1\}XW.
\end{align*}
In particular, the static insertion has population indices
\[
 (\mathcal T_{1,S})_{ac}
 =X^2\sum_b 2b_a\Theta_{ab}m_{2b}W_{bc}.
\]
The internal population index $b$ is summed. This corresponds to the matrix
product $\diag(b)\Theta\,2\diag(m_2)W$, rather than retaining $b$ as a
free source index. With this contraction its energy contribution plus the
static source vertex agrees with the earlier conditional-label adjoint
calculation: $23.37895511$ for 60:40 and $19.51649280$ for 80:20.

\subsection{Source covariance, energy and strict first-order PR}
Let $\mathcal L_\mu(\omega)=H(\omega)c_\mu(\omega)$ and define
\[
 \mathcal B_\mu(\omega)=2m\mathcal L_\mu(\omega)
            +\widehat T_\mu^{\rm dyn}(\omega),\qquad
 T_\mu^{\rm dyn}(\tau)=
       \E[v(t+\tau)^2\delta r(t)\mid\mu].
\]
The leading row-source cross covariances are diagonal in population:
\begin{align*}
 \operatorname{Cov}(\zeta^A,\xi)
 &=\diag\{\E_a[\mathcal B_\mu(\omega_1)
                      \mathcal L_\mu(\omega_2)]/f_a\},\\
 \operatorname{Cov}(\zeta^B,\xi)
 &=\diag\{\E_a[\mathcal L_\mu(\omega_1)
                      \mathcal B_\mu(\omega_2)]/f_a\}.
\end{align*}
Define $Y=\mathcal C\operatorname{Cov}(\zeta^A+\zeta^B,\xi)$; elementwise,
\[
 Y_{ab}=\tfrac12 b_aT_{ab}\E_b[
        \mathcal B_\mu(\omega_1)\mathcal L_\mu(\omega_2)
       +\mathcal L_\mu(\omega_1)\mathcal B_\mu(\omega_2)].
\]
The factor $\Theta_{ab}/f_b=T_{ab}$ is fixed by the source normalization.
Through first order,
\[
 \Sigma_S=\Sigma_0+\epsilon\Sigma_1,\qquad
 \Sigma_1=\diag(Q_{1a}/f_a)+Y+Y^\dagger.
\]
For $R_0=(I-\mathcal T_0)^{-1}$ and $g=f^\top R_0$, the energy coefficients
at a frequency pair are
\begin{align*}
 \mathcal E_0(\omega_1,\omega_2)&=g\Sigma_0g^\dagger,\\
 \mathcal E_1(\omega_1,\omega_2)&=g\Sigma_1g^\dagger
       +2\operatorname{Re}\{g\mathcal T_1R_0\Sigma_0g^\dagger\}.
\end{align*}
Integrate both frequencies to obtain $\mathcal E_0,\mathcal E_1$. In particular,
the source vertex is dressed by $R_0$ on both sides; a bare paired-ladder source
without the corresponding Schur-complement dressing omits some insertions.

The following expansion concerns the deterministic ensemble-saddle pieces,
denoted by a subscript $\mathrm{det}$, rather than a prediction for each
quenched network. Keeping the numerator to the same order gives
\begin{align*}
 \nu_{\mathrm{det}}^2&=\nu_0^2+2\epsilon\nu_0\nu_1+O(\epsilon^2),\\
 \mathcal E_{\mathrm{det}}^{-1}&=\mathcal E_0^{-1}
             [1-\epsilon\mathcal E_1/\mathcal E_0]+O(\epsilon^2),\\
 d_{\mathrm{det}}&=d_0+\epsilon d_1+O(\epsilon^2),\qquad
 d_1=d_0\left(2\nu_1/\nu_0-\mathcal E_1/\mathcal E_0\right).
\end{align*}
If a centered, tight $O_p(\epsilon)$ quenched fluctuation expansion also holds,
an individual network has an additional first-order term
\[
 d_N=d_{\mathrm{det}}+\epsilon d_0
   (2Z_{\nu,N}/\nu_0-Z_{\mathcal E,N}/\mathcal E_0)+o_p(\epsilon).
\]
Obtaining an ensemble expectation from this statement needs moment control;
neither the full fluctuation law nor its variance is proved by the present
closure. Label sampling gives only a component model for these fluctuations.
These are formal remainder orders under the stated closure assumptions, not
established uniform bounds on finite simulated networks. Evaluating a ratio
of linearly shifted numerator and denominator adds selected second-order
powers and is not the strict first-order expression.

The corrected coefficient calculation gives
\begin{center}\small
\begin{tabular}{lrrrrr}\toprule
composition & $\nu_1/\nu_0$ & $\mathcal E_1/\mathcal E_0$ &
 $d_1/d_0$ & $d_0$ & $d_0+d_1/\sqrt{1000}$\\\midrule
60:40&2.78467&4.53503&1.03431&.02989352&.03087128\\
80:20&3.19158&4.57820&1.80495&.03432623&.03628548\\\bottomrule
\end{tabular}
\end{center}
The corresponding relative PR shifts are $+3.27\%$ and $+5.71\%$. The
energy coefficient at 60:40 is
\[
 -24.00948+16.68511+3.12127+6.69385+1.05686=3.54761.
\]
The terms are, respectively, the local law, static response insertion,
dynamic trajectory-weight insertion, static source vertex, and dynamic
source vertex. This cancellation explains why the final small number alone
cannot establish numerical control at $N=1000$.

\subsection{Coherent feedback and the exact mean residual}
The nonzero connectivity mean is not a perturbative centered-disorder vertex.
For normalized population indicators $V$, write $J_R=J_b+VBV^\top$ and
let $S_0=D_1H$, $M_b=(I-S_0J_b)^{-1}$. The exact causal response identity is
\[
 M=M_b+M_bS_0VB(I-V^\top M_bS_0VB)^{-1}V^\top M_b.
\]
For time-dependent gains the small inverse is an inverse of time operators.
Eliminating this population block requires an appropriate Schur complement;
setting $B=0$ before the calculation does not derive the screened bulk law.
If a projected bulk response has a uniformly bounded inverse on a specified
frequency range and the limiting $B/\sqrt N$ is nonsingular, the direct and
indirect coherent mean probe couplings cancel at order $N^{-1/2}$ on that
range. This is not uniform high-frequency or causal $L^2$ operator counting.

Exact stationarity also retains a finite-network mean residual. Define
\[
 \eta_i=\sum_jX_{ij}m_j,\quad
 \bar\eta_a=N_a^{-1}\sum_{i\in a}\eta_i,\quad
 \bar\mu_a=N_a^{-1}\sum_{i\in a}\langle x_i\rangle_t,
 \quad m^*=-(J^0F)^{-1}I.
\]
Averaging the state equation gives
\begin{align*}
 \bar\mu&=\sqrt N J^0F(\bar m-m^*)+h^{\rm ext}+\bar\eta,\\
 \boxed{\bar m-m^*=\epsilon(J^0F)^{-1}
             (\bar\mu-h^{\rm ext}-\bar\eta).}
\end{align*}
The residual $\bar\eta$ cannot be dropped merely because each centered edge
has zero ensemble mean: $m_j$ depends on the network. If
$\bar\mu=h_0+\epsilon h_1+\cdots$,
$\bar\eta=\epsilon\eta_1+\cdots$, and $h^{\rm ext}=h_0$, the deterministic
mean-rate correction begins as
$\epsilon^2(J^0F)^{-1}(h_1-\eta_1)$. This explains the zero first-order
constraint above, but does not fix all finite-loop corrections.
Keeping an $\epsilon(J^0F)^{-1}h_1$ term inside the coefficient equations is
a partial finite-gain resummation only after specifying what $h_1$ excludes
and how the label mean and $\eta_1$ are represented. It is not an exact
finite-network balance equation without the residual term.

\subsection{Formal power counting and the stationary-limit assumption}
At a finite time grid, introduce the dynamical response fields enforcing the
state equations and let $A_{ij}$ denote their time/replica contraction with
the rate of cell $j$. Independence of the centered edges gives an expansion
of the disorder-averaged generating functional whose $n$th-cumulant term is
proportional to
\[
 \sum_{ij}\frac{\kappa_n(X_{ij})}{n!}A_{ij}^{n}.
\]
Dense weak-coupling counting gives $N^2N^{-n/2}=N^{2-n/2}$. Thus the second
cumulant contributes an order-$N$ bulk saddle action, the third an
order-$\sqrt N$ perturbation, and the fourth an order-one term. About a
nonsingular saddle, the third cumulant shifts both the saddle and the
quadratic fluctuation operator at relative order $\epsilon$. In particular,
the latter fluctuation sector contains the overlap and the two triple
overlaps retained above.

For the baseline Gaussian fluctuation law, its cubic fluctuation interaction
is odd. A single insertion into a quadratic overlap observable has zero
Gaussian expectation, since the resulting integrand has odd total degree.
Ordinary fluctuation-loop corrections therefore start one order later in
this formal expansion. This argument assumes centered Gaussian fluctuations,
a stable saddle, and a consistent expansion of the observable and source;
it does not apply to a mixture of saddles or a soft/localized mode without
additional analysis. The coherent mean has a stronger order-$N^{3/2}$ action
and must be retained or eliminated before applying this centered bulk count.

The replica identity and Gaussian Price derivatives explain why the adopted
stationary calculation uses separate conditional mean gains across replicas.
They do not establish uniform bounds on arbitrary products of pathwise
tangent matrices. Passing from a fixed-grid generating functional to the
stationary spectral result still requires mixing, self-averaging and
interchange of limits. The closure makes these assumptions explicit rather
than obtaining them from the overlap identity alone.

\subsection{A positive augmented noise and a known second-order source sector}
Exact inversion of $I-\mathcal T_0-\epsilon\mathcal T_1$ repeats known
response insertions, but its strictly first-order source omits a definite
quadratic term. Let $\zeta=\zeta^A+\zeta^B$ and
$\eta=\mathcal C\zeta$. Then
\[
 \operatorname{Cov}(\xi+\epsilon\eta)
 =\Sigma_\xi+\epsilon(Y+Y^\dagger)
                 +\epsilon^2\Sigma_\eta.
\]
The last term is nonnegative as a matrix. The leading joint source is a
covariance of the three two-replica features
$\delta r^{(1)}\delta r^{(2)}$,
$s^{(1)}v^{(2)}$, and $v^{(1)}s^{(2)}$; it is positive semidefinite by
construction.

To evaluate its missing blocks, define the one-site spectra
\[
 P_{vv}=|H|^2c_\mu,\quad S_{sv}=\mathcal B_\mu H^*,\quad
 S_{vs}=S_{sv}^*,\quad S_{ss}=\widehat{\operatorname{Cov}(s(t+\tau),s(t))}.
\]
The noise blocks are diagonal in population, with entries
\begin{align*}
 (\Sigma_{AA})_{aa}&=f_a^{-1}\E_a[S_{ss}(\omega_1)P_{vv}(\omega_2)],\\
 (\Sigma_{BB})_{aa}&=f_a^{-1}\E_a[P_{vv}(\omega_1)S_{ss}(\omega_2)],\\
 (\Sigma_{AB})_{aa}&=f_a^{-1}\E_a[S_{sv}(\omega_1)S_{vs}(\omega_2)],
 \qquad \Sigma_{BA}=\Sigma_{AB}^{\dagger},\\
 \Sigma_\eta&=\mathcal C
  (\Sigma_{AA}+\Sigma_{BB}+\Sigma_{AB}+\Sigma_{BA})\mathcal C^\dagger.
\end{align*}
For $T_v(\tau)=\E[v(t+\tau)^2v(t)\mid\mu]$, the required fourth functional
is explicit:
\begin{align*}
 \operatorname{Cov}(s(t+\tau),s(t)\mid\mu)
 &=4m^2P_\mu(\tau)+2m[T_v(\tau)+T_v(-\tau)]
       +2P_\mu(\tau)^2\\
 &\quad+\operatorname{cum}(v(t+\tau),v(t+\tau),v(t),v(t)\mid\mu).
\end{align*}
This can be evaluated from a four-time Gaussian rate expectation followed
by leak convolution. It is an additional one-site functional, not an unknown
nonlinear reservoir response.

On resolved spectral support $\Omega_{\rm res}$ where $\Sigma_0$ is positive,
the joint leading
covariance of $\xi$ and $\eta$ has Schur complement
\[
 \Sigma_\eta-Y\Sigma_0^{-1}Y^\dagger\succeq0.
\]
Use the pseudoinverse with the corresponding range condition if a true null
space is present. If $R_\epsilon=(I-\mathcal T_0-\epsilon\mathcal T_1)^{-1}$,
the omitted known source contribution consequently obeys, within this
leading row-noise model,
\[
 \Delta\mathcal E_{\rm source}
 \ge\epsilon^2\iint_{\Omega_{\rm res}} f^\top R_\epsilon
       Y\Sigma_0^{-1}Y^\dagger R_\epsilon^\dagger f
       \frac{\dd\omega_1\dd\omega_2}{(2\pi)^2}.
\]
On resolved spectral support the numerical lower contributions at $N=1000$
are $.01579386$ for 60:40 and $.01256055$ for 80:20. They are $1.866\%$ and
$1.803\%$ of the respective partial-resummed energies. Thresholds between
$10^{-8}$ and $10^{-12}$ relative to the leading source scale change these
figures negligibly. The signed first-order source has small negative tail
eigenvalues of order $10^{-14}$; their integrated negative-energy fractions
are below $3\times10^{-11}$, so numerical tail negativity is not the
material effect measured by this Schur bound.

This inequality concerns an omitted source covariance of the specified
augmented model. It is not a lower bound on the true network's total
second-order energy correction: unknown second-order one-site shifts,
vertices, response corrections and fluctuation loops can have either sign.
Restoring the full augmented covariance is preferable to an arbitrary
eigenvalue clipping rule, but still does not complete that second-order theory.

\subsection{Scope of the adopted result}
The corrected closure supplies every function in its formal first-order
coefficient through one-site Gaussian integrals, leak convolutions, a linear
marginal problem and two-by-two population matrix inverses. In particular it
does not leave a stationary nonlinear-reservoir average unevaluated within
that adopted approximation. Its new trajectory weight and correctly dressed
source are meaningful analytical additions to the previous partial candidate.

The remaining claims are narrower than a full finite-network solution.
Uniform stationary and delocalization bounds are not proved; the single-synapse
EE amplitude is about $.55$ at $N=1000$; large sectors cancel; known omitted
source covariance already contributes at the percent level; and the complete
quenched law, including correlations beyond a sampled static-label model,
is not determined. An influence-function calculation for independent labels
is a component model for fluctuations, not a rigorous lower bound on their
total variance. The leading theory and this formal extension must be tested
against independently sampled networks, with active-state selection and
finite-record estimation described separately.
